\section{Yêu cầu về hệ thống}

\subsection{Yêu cầu chức năng}

Dựa trên phân tích các vấn đề cốt lõi trong công tác quản lý Đoàn -- Hội tại Trường Đại học Bách Khoa, hệ thống cần đáp ứng các yêu cầu chức năng sau:

\subsubsection{Quản lý Người dùng và Tổ chức}

\begin{itemize}
    \item \textbf{Quản lý hồ sơ điện tử:}
    \begin{itemize}
        \item Hệ thống phải lưu trữ và quản lý thông tin định danh của Đoàn viên và Hội viên bao gồm: mã số sinh viên (MSSV), họ tên, email, số điện thoại, ngày sinh, giới tính, tình trạng Đoàn viên/Hội viên.
        \item Hệ thống phải lưu trữ thông tin học vụ: khoa, bộ môn, lớp, niên khóa, chương trình đào tạo.
        \item Hệ thống phải tự động tích lũy điểm CTXH khi sinh viên hoàn thành hoạt động và được chấm điểm.
        \item Hệ thống phải hỗ trợ import hàng loạt từ file Excel (18 cột) với khả năng tự động tạo tổ chức/khoa/lớp nếu chưa tồn tại trong hệ thống.
    \end{itemize}
    
    \item \textbf{Quản lý cơ cấu tổ chức:}
    \begin{itemize}
        \item Hệ thống phải hỗ trợ phân cấp tổ chức theo mô hình: Trường $\rightarrow$ Khoa (Đoàn/Hội cơ sở) $\rightarrow$ Lớp (Chi đoàn/Chi hội).
        \item Hệ thống phải hỗ trợ 8 loại tổ chức: Đoàn trường, Đoàn cơ sở, Chi đoàn, Hội sinh viên trường, Liên chi hội, Chi hội, CLB trường, CLB khoa.
        \item Hệ thống phải cho phép quản lý khoa (Department) và lớp (Class) độc lập, đồng thời liên kết với tổ chức tương ứng.
        \item Hệ thống phải cho phép gán sinh viên vào nhiều tổ chức với các vai trò: member, admin, hoặc super\_admin.
    \end{itemize}
    
    \item \textbf{Phân quyền truy cập (RBAC):}
    \begin{itemize}
        \item Hệ thống phải áp dụng mô hình Role-Based Access Control với context theo tổ chức.
        \item Hệ thống phải phân chia quyền theo 2 cấp độ: System (toàn hệ thống) và Organization (theo đơn vị cụ thể).
        \item Hệ thống phải đảm bảo cán bộ chỉ được truy cập và chỉnh sửa dữ liệu thuộc phạm vi quản lý (organization scope).
        \item Hệ thống phải hỗ trợ gán nhiều vai trò cho cùng một user ở các tổ chức khác nhau.
        \item Hệ thống phải cung cấp các vai trò mặc định: Super Admin, Quản trị viên tổ chức, Quản lý, Biên tập viên, Người xem.
    \end{itemize}
\end{itemize}

\subsubsection{Quản lý Hoạt động CTXH}

\begin{itemize}
    \item \textbf{Tạo và quản lý hoạt động:}
    \begin{itemize}
        \item Hệ thống phải cho phép cán bộ tạo, chỉnh sửa, xóa mềm (soft delete) thông tin hoạt động bao gồm: tên, mô tả, thời gian bắt đầu/kết thúc, địa điểm, số lượng tham gia, tổng số ngày CTXH.
        \item Hệ thống phải quản lý 4 trạng thái hoạt động: DRAFT (Nháp), ACTIVE (Đang hoạt động), COMPLETED (Đã kết thúc), CANCELLED (Đã hủy).
        \item Hệ thống phải kiểm tra ràng buộc: chỉ cho phép thêm sinh viên khi hoạt động ở trạng thái ACTIVE; chỉ cho phép chấm điểm khi hoạt động ở trạng thái ACTIVE hoặc COMPLETED.
        \item Hệ thống phải hỗ trợ phân loại hoạt động theo danh mục (Category) với cấu trúc phân cấp cây.
    \end{itemize}
    
    \item \textbf{Quản lý danh sách tham gia:}
    \begin{itemize}
        \item Hệ thống phải cho phép cán bộ thêm sinh viên vào hoạt động theo 2 cách:
        \begin{enumerate}
            \item Thêm thủ công từng sinh viên bằng cách nhập mã số sinh viên.
            \item Upload file Excel chứa danh sách sinh viên (hệ thống tự động parse và validate từng dòng).
        \end{enumerate}
        \item Hệ thống phải theo dõi 2 trạng thái tham gia: PENDING (Chưa chấm điểm), GRADED (Đã chấm điểm).
        \item Hệ thống phải ghi nhận các thông tin: ngày đăng ký tham gia, ngày chấm điểm, số ngày CTXH thực tế, ghi chú, mã người chấm điểm.
    \end{itemize}
    
    \item \textbf{Chấm điểm CTXH:}
    \begin{itemize}
        \item Hệ thống phải cho phép cán bộ chấm điểm thủ công bằng cách nhập số ngày thực tế (actual\_days) và ghi chú cho từng sinh viên.
        \item Hệ thống phải hỗ trợ upload file Excel để chấm điểm hàng loạt, kèm validation: số ngày không được vượt quá total\_days của hoạt động, không cho phép chấm trùng lặp.
        \item Hệ thống phải cung cấp chức năng tải template Excel chứa danh sách sinh viên đã tham gia để cán bộ điền điểm.
        \item Hệ thống phải tự động cộng dồn điểm CTXH vào tổng điểm trong hồ sơ sinh viên khi chấm điểm thành công.
    \end{itemize}
    
    \item \textbf{Thống kê và báo cáo:}
    \begin{itemize}
        \item Báo cáo số lượng hoạt động theo danh mục, tổ chức, trạng thái.
        \item Thống kê tỷ lệ hoàn thành của sinh viên theo hoạt động.
        \item Xuất báo cáo điểm CTXH cá nhân và tổng hợp của hoạt động (file Excel).
    \end{itemize}
\end{itemize}

\subsubsection{Quản lý Đóng Phí}

\begin{itemize}
    \item \textbf{Tạo đợt đóng phí:}
    \begin{itemize}
        \item Cán bộ tạo Payment Campaign với thông tin: tên đợt, mô tả, thời gian bắt đầu/kết thúc, tổ chức phụ trách.
        \item Quản lý trạng thái đợt: DRAFT, ACTIVE, COMPLETED, CANCELLED.
        \item Cấu hình mức phí (Fee Level) theo đối tượng: sinh viên, cán bộ, hoặc custom với số tiền khác nhau.
    \end{itemize}
    
    \item \textbf{Phân công đóng phí:}
    \begin{itemize}
        \item Cán bộ phân công sinh viên vào đợt đóng phí (Payment Assignment) với mức phí tương ứng.
        \item Hỗ trợ phân công hàng loạt theo lớp, khoa, hoặc danh sách MSSV.
        \item Theo dõi trạng thái thanh toán: PENDING, PAID, CANCELLED.
        \item Ghi nhận thông tin: hạn đóng, số tiền, ghi chú.
    \end{itemize}
    
    \item \textbf{Thanh toán trực tuyến qua PayOS:}
    \begin{itemize}
        \item Sinh viên xem danh sách các khoản phí được phân công trong trang "Đóng phí của tôi".
        \item Click "Thanh toán ngay" để chuyển đến cổng PayOS.
        \item Hỗ trợ các phương thức: chuyển khoản ngân hàng, ví điện tử (VNPay, Momo, ZaloPay).
        \item Webhook tự động cập nhật trạng thái thanh toán sau khi giao dịch thành công.
    \end{itemize}
    
    \item \textbf{Cấu hình PayOS theo từng đợt:}
    \begin{itemize}
        \item Mỗi đợt đóng phí có thể cấu hình riêng: Client ID, API Key, Checksum Key.
        \item Hỗ trợ kế thừa cấu hình từ tổ chức cha nếu không tự cấu hình.
        \item Lưu lịch sử giao dịch (Payment History): mã giao dịch, số tiền, trạng thái, thời gian.
    \end{itemize}
    
    \item \textbf{Báo cáo thu phí:}
    \begin{itemize}
        \item Theo dõi tiến độ thu phí: tổng số người được phân công, số người đã đóng, tỷ lệ hoàn thành.
        \item Thống kê doanh thu theo đợt, theo tổ chức, theo thời gian.
        \item Xuất báo cáo Excel danh sách sinh viên đã/chưa đóng phí.
    \end{itemize}
\end{itemize}

\subsubsection{Quản lý Biểu mẫu và Chứng từ}

\begin{itemize}
    \item \textbf{Upload và parse template:}
    \begin{itemize}
        \item Quản trị viên upload file Word (.docx) chứa các placeholder dạng \texttt{\{\{field\_name\}\}}.
        \item Hệ thống tự động parse file bằng thư viện Docxtemplater, trích xuất danh sách merge fields.
        \item Upload file template lên S3 (Long Vân), lưu URL vào Database.
    \end{itemize}
    
    \item \textbf{Cấu hình fields:}
    \begin{itemize}
        \item Quản trị viên cấu hình từng field: tên hiển thị, kiểu dữ liệu (TEXT, NUMBER, DATE, SELECT...), required, validation rules.
        \item Hỗ trợ các field đặc biệt: DATE\_PARTS (tách ngày/tháng/năm), SELECT với options tùy chỉnh.
        \item Quản lý trạng thái template: PROCESSING (đang xử lý), ACTIVE (sẵn sàng sử dụng), INACTIVE (vô hiệu hóa).
    \end{itemize}
    
    \item \textbf{Điền và nộp biểu mẫu:}
    \begin{itemize}
        \item Sinh viên chọn template, hệ thống hiển thị form động dựa trên cấu hình fields.
        \item Validate dữ liệu trên client (ReactJS) và server (NestJS) trước khi submit.
    \end{itemize}
    
    \item \textbf{Generate document:}
    \begin{itemize}
        \item Khi submit, hệ thống download template từ S3, load vào Docxtemplater.
        \item Transform dữ liệu: chuyển DATE thành DD/MM/YYYY, tạo field\_day/field\_month/field\_year cho DATE\_PARTS.
        \item Render document với dữ liệu đã transform, generate buffer.
        \item Upload file output lên S3 với đường dẫn: \texttt{submissions/\{userId\}/\{timestamp\}\_\{fileName\}.docx}.
        \item Cập nhật submission: status = COMPLETED, outputFileUrl, completedAt.
    \end{itemize}
    
    \item \textbf{Lưu trữ và tra cứu:}
    \begin{itemize}
        \item Lưu metadata submission trong Database: templateId, userId, formData (JSON), status, timestamps.
        \item Sinh viên xem danh sách biểu mẫu đã nộp, tải lại file đã generate.
        \item Hỗ trợ preview trực tuyến qua Office Online Viewer (presigned URL từ S3).
    \end{itemize}
\end{itemize}

\subsubsection{Chatbot hỏi đáp hoạt động công tác xã hội}

\begin{itemize}
    \item Trả lời câu hỏi về thông tin hoạt động công tác xã hội hiện có trong hệ thống
\end{itemize}

\subsection{Yêu cầu phi chức năng}

Hệ thống được thiết kế với các yêu cầu phi chức năng sau đây, dựa trên công nghệ đã triển khai và nhu cầu thực tế của đối tượng sử dụng:

\subsubsection{Bảo mật và An toàn thông tin}

\begin{itemize}
    \item \textbf{Xác thực và phân quyền:}
    \begin{itemize}
        \item Hệ thống phải mã hóa mật khẩu bằng thuật toán bcrypt với salt rounds = 10 trước khi lưu vào cơ sở dữ liệu.
        \item Hệ thống phải xác thực người dùng qua JWT với thời gian hết hạn access token là 15 phút và refresh token là 7 ngày.
        \item Hệ thống phải phân quyền theo mô hình RBAC với context tổ chức, đảm bảo mỗi user chỉ được truy cập dữ liệu thuộc phạm vi được gán.
        \item Hệ thống phải kiểm tra quyền trước khi cho phép thực thi mọi API liên quan đến tạo, sửa, xóa dữ liệu quan trọng.
    \end{itemize}
    
    \item \textbf{Bảo vệ dữ liệu:}
    \begin{itemize}
        \item Hệ thống phải validate và sanitize mọi input từ người dùng trên cả client và server để phòng chống SQL Injection, XSS.
        \item Hệ thống phải sử dụng ORM (TypeORM) để xây dựng query, tránh concatenate SQL string trực tiếp.
        \item Hệ thống phải áp dụng soft delete cho các entity quan trọng (users, activities, organizations) để có thể khôi phục khi cần thiết.
        \item Hệ thống phải lưu trữ timestamp (createdAt, updatedAt, deletedAt) và thông tin người thao tác cho mọi entity quan trọng phục vụ audit trail.
    \end{itemize}
    
    \item \textbf{Triển khai bảo mật:}
    \begin{itemize}
        \item Hệ thống phải sử dụng HTTPS với chứng chỉ SSL/TLS hợp lệ trên môi trường production.
        \item Hệ thống phải cấu hình CORS chỉ cho phép request từ domain chính thức của ứng dụng.
        \item Hệ thống phải kiểm tra extension và MIME type khi upload file, chỉ chấp nhận .xlsx và .docx với kích thước tối đa 10MB.
    \end{itemize}
\end{itemize}

\subsubsection{Hiệu năng}

\begin{itemize}
    \item \textbf{Tối ưu hóa cơ sở dữ liệu:}
    \begin{itemize}
        \item Hệ thống phải đánh index trên các trường tra cứu thường xuyên (code, email, student\_code) để đảm bảo thời gian truy vấn $<$ 100ms với bảng có $<$ 100,000 bản ghi.
        \item Hệ thống phải áp dụng pagination cho mọi danh sách với kích thước trang mặc định 10-50 bản ghi.
        \item Hệ thống phải sử dụng eager/lazy loading hợp lý để tránh N+1 query problem, đảm bảo số lượng query $\leq$ 10 cho mỗi request.
        \item Hệ thống phải xử lý batch operations (import Excel $>$ 1000 dòng) theo chunk để tránh timeout.
    \end{itemize}
    
    \item \textbf{Xử lý file:}
    \begin{itemize}
        \item Hệ thống phải upload file lên S3 storage thay vì lưu trực tiếp vào Database để tránh quá tải.
        \item Hệ thống phải parse file Excel theo chunk (200 dòng/lần) để xử lý file lớn mà không quá tải bộ nhớ.
        \item Hệ thống phải generate document Word bất đồng bộ và cập nhật trạng thái khi hoàn thành, timeout tối đa 60 giây.
    \end{itemize}
    
    \item \textbf{Tối ưu Frontend:}
    \begin{itemize}
        \item Hệ thống phải sử dụng client-side caching để giảm số lần gọi API lặp lại cho cùng dữ liệu (cache time: 5 phút).
        \item Hệ thống phải lazy load các module và route để đảm bảo initial bundle size $<$ 500KB (gzipped).
        \item Hệ thống phải tối ưu re-render để đảm bảo thời gian phản hồi UI $<$ 100ms khi người dùng tương tác.
    \end{itemize}
\end{itemize}

\subsubsection{Khả năng triển khai và Mở rộng}

\begin{itemize}
    \item \textbf{Containerization:}
    \begin{itemize}
        \item Đóng gói ứng dụng bằng Docker: backend (NestJS), frontend (ReactJS + Nginx), database (PostgreSQL).
        \item Sử dụng Docker Compose để orchestrate các container trong môi trường development và production.
        \item Multi-stage build trong Dockerfile để giảm kích thước image production.
    \end{itemize}
    
    \item \textbf{Kiến trúc modular:}
    \begin{itemize}
        \item Tách biệt các module: Users, Organizations, CTXH, Payment, Templates, Auth.
        \item Mỗi module có structure riêng: entities, services, controllers, DTOs.
        \item Dependency injection với NestJS để dễ dàng test và thay thế implementation.
    \end{itemize}
    
    \item \textbf{Lưu trữ phân tán:}
    \begin{itemize}
        \item File storage sử dụng S3-compatible (Long Vân), tách biệt khỏi database server.
        \item Presigned URL cho file private, cho phép download trực tiếp từ S3 mà không qua backend.
    \end{itemize}
\end{itemize}

\subsubsection{Giám sát và Logging}

\begin{itemize}
    \item \textbf{Application monitoring:}
    \begin{itemize}
        \item Tích hợp Prometheus để thu thập metrics: request count, response time, error rate.
        \item Visualize metrics với Grafana dashboard, theo dõi real-time.
        \item Cấu hình alerts khi phát hiện bất thường (error rate tăng cao, memory leak).
    \end{itemize}
    
    \item \textbf{Logging:}
    \begin{itemize}
        \item Sử dụng NestJS Logger để ghi log với các level: error, warn, info, debug.
        \item Log các thao tác quan trọng: login, create/update/delete user, chấm điểm CTXH, thanh toán.
        \item Lưu trữ logs để phục vụ debugging và audit.
    \end{itemize}
    
    \item \textbf{Error handling:}
    \begin{itemize}
        \item Global exception filter để bắt và format lỗi thống nhất.
        \item Trả về error message rõ ràng cho frontend hiển thị (không expose stack trace ở production).
        \item Transaction rollback tự động khi có lỗi trong các thao tác liên quan nhiều bảng.
    \end{itemize}
\end{itemize}

\subsubsection{Giao diện và Trải nghiệm người dùng}

\begin{itemize}
    \item \textbf{Responsive design:}
    \begin{itemize}
        \item Sử dụng Tailwind CSS với breakpoints để hỗ trợ desktop, tablet, mobile.
        \item Test giao diện trên các kích thước màn hình phổ biến (1920x1080, 1366x768, 375x667).
    \end{itemize}
    
    \item \textbf{Dark mode:}
    \begin{itemize}
        \item Hỗ trợ chế độ sáng/tối, tự động lưu preference của user.
        \item Áp dụng dark mode classes cho tất cả component với Tailwind CSS.
    \end{itemize}
    
    \item \textbf{Đa ngôn ngữ:}
    \begin{itemize}
        \item Tích hợp i18next để hỗ trợ tiếng Việt và tiếng Anh.
        \item Tách file translation theo module để dễ quản lý và mở rộng.
    \end{itemize}
    
    \item \textbf{UX improvements:}
    \begin{itemize}
        \item Loading states cho các thao tác mất thời gian (upload file, generate document).
        \item Toast notifications cho feedback tức thì khi thành công/lỗi.
        \item Confirmation dialog cho các thao tác nguy hiểm (xóa user, hủy hoạt động).
    \end{itemize}
\end{itemize}

\subsubsection{Code Quality và Maintainability}

\begin{itemize}
    \item \textbf{Coding standards:}
    \begin{itemize}
        \item Backend: tuân thủ NestJS style guide, sử dụng TypeScript strict mode.
        \item Frontend: ESLint + Prettier để format code tự động, maintain consistency.
        \item Naming convention rõ ràng: camelCase cho biến/hàm, PascalCase cho class/component.
    \end{itemize}
    
    \item \textbf{Documentation:}
    \begin{itemize}
        \item API documentation tự động với Swagger UI (\texttt{/api-docs}).
        \item Inline comments cho logic phức tạp trong code.
        \item README files cho từng module, giải thích cách sử dụng và dependencies.
    \end{itemize}
    
    \item \textbf{Version control:}
    \begin{itemize}
        \item Sử dụng Git với branch strategy: main (production), develop (staging), feature branches.
        \item Commit message theo format: type(scope): description (VD: feat(ctxh): add grading from Excel).
    \end{itemize}
\end{itemize}
